% Nouveau EGA V (cote 25): archive cote 25, site page 76 (PDF page 77).
% typescript. Draft compared: none.
% AI-transcribed by Claude Opus 5.5 (claude-opus-5-5, adaptive thinking, effort high), Message Batches API, batch ega25b (msgbatch_01A7n5YuXubUPNK1zsepi8RG), 2026-11-13; not yet checked by eye.
\scanpage{25/76}
\mspage{71}

\unsure{K|k} soit parfait assure que X' est lisse sur k (car X' est en tous cas
régulier), et il assure également que $Z = X' - X$ , muni de la structure
réduite induite, est étale sur k. Traitons maintenant le cas général,
dans lequel il est loisible de supposer $n \geqslant 2$.

\xout{Comme X est affine donc quasi-projectif,} On peut évidemment suppo-%
ser X \xout{\unsure{lisse et}} affine donc quasi-projectif, puis, quitte à remplacer X
par une clôture projective, on peut supposer X projectif, sous réserve
toutesfois de prouver que \uline{tout} voisinage de x contient un voisinage
ouvert U admettant un morphisme élémentaire $U \to V$. De même, remplaçant
X par son normalisé (fini sur X donc projectif, réf.), ce qui ne le
change pas au voisinage de x, on peut supposer X normal, donc k étant
parfait, géométriquement normal sur k. \xout{Choisissons une immersion} Le
bénéfice de cette hypothèse est que l'ensemble Z des points de X en
lesquels X n'est pas lisse sur k est de codim $\geqslant 2$. \xout{\unsure{Soit Noter qui}}
Choisissons une immersion projective $i : X \longrightarrow P^r$ , \xout{\unsure{On voit on peut}}
\xout{\unsure{supposer que X n'est contenu dans Alors}} Comme on obtient un système
fondamental de voisinages de x dans $\underline{P}^{\hand{r}}$ en prenant des sections ne
s'annulant pas en x des divers $\underline{O}_P(n)$ , $n > 0$, on en conclut que tout
voisinage de x contient un voisinage de la forme X-Y, où Y est une
partie fermée de X , contenant Z, purement de dimension n-1, et telle que
$x \notin Y$. Nous munirons Y de la structure réduite induite, de sorte que
(k étant parfait) l'ensemble singulier de Y est de dimension $\leqslant$ n-2. \margin{\unsure{$\nearrow$}}
Quitte à agrandir l'ensemble Z précédent, on trouve un fermé $Z \subset Y$, de
dimension $\leqslant$ n-2, contenant l'ensemble géom. singulier de X \uline{et} de Y.

L'idée de la démonstration est de fibrer X par ses intersections
avec les sous-variétés linéaires L de P , de codimension n-1, contenant

\mspage{76}
